\documentclass[11pt]{article}

\usepackage{amsmath}
\usepackage{amssymb}
\usepackage{amsthm}
\usepackage[margin=1in]{geometry}

\newtheorem{definition}{Definition}[section]
\newtheorem{theorem}{Theorem}[section]
\newtheorem{example}{Example}[section]
\newtheorem*{remark}{Remark}

\title{Basic Mathematics --- Chapter 1 Notes}
\author{Zach Lysobey}
\date{\today}

\begin{document}

\maketitle

\section{Integers}

\begin{definition}
  An \emph{integer} is a whole number, positive, negative, or zero:
  $\ldots, -2, -1, 0, 1, 2, \ldots$
\end{definition}

\begin{definition}
  The integers $1, 2, 3, \ldots$ are called the \emph{positive integers}.
\end{definition}

\begin{definition}
  An integer that is either $0$ or positive is called a \emph{natural
  number}. Thus the natural numbers are $0, 1, 2, 3, \ldots$
\end{definition}

\begin{remark}
  Note Lang's convention: $0$ \emph{is} a natural number. Some books use
  ``natural numbers'' to mean only the positive integers $1, 2, 3, \ldots$
\end{remark}

\begin{definition}
  When the integers are pictured as points on a line, the point
  corresponding to $0$ is called the \emph{origin}. Positive integers lie
  to the right of the origin, negative integers to the left.
\end{definition}

\subsection*{First rules for addition}

On the number line, adding a positive integer means moving that many
units to the right; adding a negative integer means moving to the left.
Two rules are taken as assumptions:

\begin{description}
  \item[N1.] $a + 0 = a$ \quad for any integer $a$.
  \item[N2.] $a + (-a) = 0$ \quad and also \quad $-a + a = 0$.
\end{description}

Geometrically, N2 says that $a$ and $-a$ lie on opposite sides of the
origin, at the same distance.

\begin{definition}
  Because of the property $a + (-a) = 0$ (rule N2), the number $-a$ is
  called the \emph{additive inverse} of $a$: it is the number which,
  added to $a$, gives $0$.
\end{definition}

\subsection*{``Minus'' vs.\ ``negative''}

For any number $a$, the number $-a$ should be read ``minus $a$,'' not
``negative $a$.'' The reason: $-a$ is not necessarily negative. If $a$ is
itself negative, then $-a$ is positive; for instance, if $a = -3$, then
$-a = -(-3) = 3$. So ``negative $a$'' wrongly suggests that $-a$ is a
negative number, while ``minus $a$'' just names the operation of taking
the additive inverse. The word \emph{negative} is reserved for numbers
less than $0$.

\section{Rules for Addition}

Addition of integers obeys two basic rules. Like N1 and N2, these are
assumed, not proved; other rules are then proved from them.

\begin{description}
  \item[Commutativity.] If $a$, $b$ are integers, then
    \[ a + b = b + a. \]
  \item[Associativity.] If $a$, $b$, $c$ are integers, then
    \[ (a + b) + c = a + (b + c). \]
\end{description}

Because of associativity, parentheses are unnecessary in a sum of several
terms: we simply write $a + b + c$. By applying associativity and
commutativity repeatedly, a sum of integers can be taken in any order.

As a matter of convention, we write
\[ a - b \quad \text{instead of} \quad a + (-b). \]
Subtraction is not a new operation: it is addition of the additive
inverse.

From the rules above we can prove:

\begin{description}
  \item[N3.] If $a + b = 0$, then $b = -a$ and $a = -b$.
\end{description}

\begin{proof}
  Add $-a$ to both sides of $a + b = 0$. The left side is
  $-a + a + b = 0 + b = b$, and the right side is $-a + 0 = -a$, so
  $b = -a$. Similarly, adding $-b$ to both sides gives $a = -b$.
\end{proof}

In words: the additive inverse of $a$ is the \emph{only} number that
added to $a$ gives $0$.

\begin{description}
  \item[N4.] $a = -(-a)$ \quad for every integer $a$.
\end{description}

\begin{proof}
  We have $-a + a = 0$ by N2. Apply N3 to this sum (with $-a$ in the
  role of $a$, and $a$ in the role of $b$): it gives $a = -(-a)$.
\end{proof}

Note that N4 holds whether $a$ is positive, negative, or $0$. Piling up
minus signs in front of $a$ gives $a$ or $-a$ alternately:
$3 = -(-3) = -(-(-(-3)))$.

\begin{remark}[Parsing repeated plus and minus signs]
  A minus sign binds only to the single term immediately after it. Thus
  \[ 1 + 2 - 3 + 4 \quad \text{means} \quad 1 + 2 + (-3) + 4 = 4, \]
  and once every term is written as an addition, associativity lets us
  regroup freely, e.g.\ as $1 + 2 + (-3 + 4)$. It does \emph{not} mean
  \[ 1 + 2 - (3 + 4) = -4. \]
  To subtract a whole group of terms, the parentheses must be written
  explicitly --- and then the minus sign distributes over the group
  (rule N5, below).
\end{remark}

For any integers $a$, $b$ we have $-(a + b) = -a + (-b)$, or in other
words:

\begin{description}
  \item[N5.] $-(a + b) = -a - b$.
\end{description}

\begin{proof}
  By N3, to show that $-a - b$ is the additive inverse of $a + b$, it
  suffices to show that the two add to $0$. And indeed
  \begin{align*}
    (a + b) + (-a - b) &= a + b - a - b \\
                       &= a - a + b - b
                          && \text{by commutativity and associativity} \\
                       &= 0 + 0 = 0. \qedhere
  \end{align*}
\end{proof}

\begin{example}
  \begin{align*}
    -(3 + 5)  &= -3 - 5 = -8, \\
    -(-4 + 5) &= -(-4) - 5 = 4 - 5 = -1, \\
    -(3 - 7)  &= -3 - (-7) = -3 + 7 = 4.
  \end{align*}
  Be careful when taking the negative of a sum which itself involves
  negative numbers: each term's sign flips, using $-(-a) = a$ where
  needed.
\end{example}

\section{Rules for Multiplication}

\end{document}
